% !PRINTED 560
\weberpage{560}
%<<CONT
aus der Bemerkung, dass die Substitution $(r, r_h)$ unter den $\eta_k$
die Substitution $(\eta_k, \eta_{k+h})$ hervorruft, wenn der Index von $\eta$
nach dem Modul $e$ genommen wird. Die Gruppe von $F_e(x) = 0$
besteht daher aus den Potenzen der cyklischen Permutation

\[
(\eta, \eta_1, \eta_2 \ldots \eta_{e-1}).
\]

\begin{center}
\S.\,168.

\gesperrt{Die Gauss'sche Methode zur Berechnung der\\
Resolventen.}
\end{center}

Um in concreten Fällen die Auflösung der Kreistheilungsgleichung
wirklich durchzuführen, kommt es nach dem, was wir
im vorigen Paragraphen entwickelt haben, nur darauf an, die
Producte je zweier conjugirter Perioden als lineare Functionen
derselben Perioden darzustellen. Dafür hat \gesperrt{Gauss} ein einfaches
Verfahren angegeben, das wir jetzt kennen lernen wollen. Wir
müssen dazu die Bezeichnung der Perioden ein wenig verändern.

Wir schreiben zwei beliebige der Perioden \S.\,167, (12) in
der Weise

\weq{1}{\begin{aligned}
\eta^{(\lambda)} &= r^\lambda + r^{\lambda'} + r^{\lambda''} + \cdots\\
\eta^{(\mu)} &= r^\mu + r^{\mu'} + r^{\mu''} + \cdots
\end{aligned}}

so dass, wenn $\lambda \equiv g^h \ (\mathrm{mod}\ n)$ ist, $\eta^{(\lambda)} = \eta_h$ oder auch
$\eta^{(\lambda)} = \eta_{\mathrm{ind}\,\lambda}$; es ist dann

\weq{2}{\begin{aligned}
\lambda' &\equiv \lambda g^e, & \lambda'' &\equiv \lambda g^{2e} \cdots\\
\mu' &\equiv \mu g^e, & \mu'' &\equiv \mu g^{2e} \cdots
\end{aligned}\ (\mathrm{mod}\ n);}

lassen wir also die Zeichen $s$ und $t$ je ein volles Restsystem
nach dem Modul $f$ durchlaufen, so ist

\[
\eta^{(\lambda)} = \sum^{s} r^{\lambda g^{se}}, \qquad \eta^{(\mu)} = \sum^{t} r^{\mu g^{te}}.
\]

Das Product davon ist

\[
\eta^{(\lambda)} \eta^{(\mu)} = \sum^{s}\sum^{t} r^{\lambda g^{se} + \mu g^{te}}.
\]

Halten wir bei der Bildung dieser Summe zunächst $s$ fest,
und summiren in Bezug auf $t$, so dürfen wir $t$ durch $t + s$ ersetzen, weil beide gleichzeitig ein volles Restsystem nach dem
Modul $f$ durchlaufen. Demnach wird
\[
\eta^{(\lambda)}\eta^{(\mu)} = \sum^{s}\sum^{t} r^{(\lambda + \mu g^{te})g^{se}}.
\]
%CONT>>
